% Part: normal-modal-logic
% Chapter: syntax-and-semantics
% Section: introduction

\documentclass[../../../include/open-logic-section]{subfiles}

\begin{document}

\olfileid{nml}{syn}{int}

\olsection{导言}

模态逻辑研究的是\emph{模态命题}以及它们之间的后承关系。模态命题的例子如下：

\begin{enumerate}
\item  $2+2=4$ 是必然的。
\item  必然地，明天可能会下雨。
\item  如果必然地可能 $!A$，那么可能 $!A$。
\end{enumerate}

可能性和必然性并非仅有的模态：其他一元联结词也被归为模态，例如，“应当是 $!A$”、“将会是 $!A$”、“Dana 知道 $!A$”，或“Dana 相信 $!A$”。

模态逻辑最早出现在亚里士多德的《解释篇》中：他最先注意到必然性蕴涵可能性，但反之不然；可能性与必然性可以相互定义；如果 $!A \land !B$ 可能为真，那么 $!A$ 可能为真且 $!B$ 可能为真，但反之不然；并且如果 $!A \to !B$ 是必然的，那么若 $!A$ 是必然的，$!B$ 也是必然的。

对模态逻辑的第一种现代研究进路源于 C.~I. Lewis 的工作，最终汇集于 Lewis 和 Langford 的《符号逻辑》（1932）一书中。Lewis 和 Langford 不满意用实质条件句来表示蕴涵：$!A \lif !B$ 是“$!A$ 蕴涵 $!B$”的一个拙劣的替代。取而代之，他们提议将蕴涵刻画为“必然地，如果 $!A$ 则 $!B$”，记作 $!A \strictif !B$。在试图厘清不同性质的过程中，Lewis 区分了五个不同的模态系统，\Log{S1}, \ldots, \Log{S4}, \Log{S5}，其中最后两个至今仍被使用。

Lewis 和 Langford 的进路是纯粹\emph{句法}的：他们提出了合理的公理和规则，并研究用这些手段能证明什么。语义进路在很长一段时间里都难以寻觅，直到 Rudolf Carnap（卡尔纳普）在《意义与必然性》（1947）中首次尝试，使用了\emph{状态描述}的概念，即原子语句的集合（即那些在该状态描述中“为真”的原子语句）。如果 $!A$ 在所有状态描述中均为真，Carnap 就定义 $!A$ 是\emph{必然为真}的。Carnap 的进路无法处理\emph{叠置}模态，因为形如“可能必然……可能 $!A$”的语句总是归约为最内层的模态。

模态语义学的重大突破随着 Saul Kripke（克里普克）的文章《模态逻辑中的一个完备性定理》（JSL 1959）而来。Kripke 的工作基于 Leibniz（莱布尼茨）的观点，即如果一个陈述在“所有可能世界”中均为真，则它是必然为真的。然而，这个观点与 Carnap 的进路有着同样的缺陷，即一个陈述在世界 $w$（或状态描述 $s$）中的真值完全不依赖于 $w$。因此 Kripke 假设世界之间由一个\emph{可达关系} $R$ 相联系，并且形如“必然 $!A$”的陈述在世界 $w$ 中为真，当且仅当 $!A$ 在所有\emph{可从 $w$ 到达的}世界 $w'$ 中均为真。提供此类进路某一版本的语义被称为 Kripke 语义，它使得多种模态逻辑得以迅猛发展。

通过 Kripke 语义的解释，模态逻辑向我们展示了\emph{关系结构}“从内部看”是什么样子。关系结构就是一个配备了二元关系的集合（例如，按社会安全号排序的班级学生集合就是一个关系结构）。但实际上，关系结构存在于各种各样的领域中：除了世界状态之间的相对可能性，还可以有某主体的认知状态通过认知可能性相关联，或者带有状态转移的动态系统状态，等等。模态逻辑可被用来为所有这些建模：第一种情形给出普通的、真势的模态逻辑；其他情形则给出认知逻辑、动态逻辑等等。

我们聚焦于一个特定的角度，模态逻辑学家称之为“对应理论”。Kripke 最重要的早期发现之一是：可达关系~$R$ 的许多性质（无论它是传递的、对称的，等等）可以通过恰当的“模态模式”\emph{在模态语言本身中得到刻画}。例如，模态逻辑学家说，$R$ 的自反性“对应”于模式“如果必然~$!A$，则~$!A$”。我们主要探讨通过组合模式 \Ax{D}、\Ax{T}、\Ax{B}、\Ax{4} 和~\Ax{5} 所得到的若干经典模态逻辑系统（例如 \Log{S4} 和 \Log{S5}）的对应理论。

\end{document}